\section{Introduction}\label{sec:intro}

Undamped or lightly damped structures keep ringing long after an impulsive
disturbance.
Passive remedies such as tuned mass dampers~\citep{denhartog1956mechanical} and
nonlinear energy sinks~\citep{vakakis2008nonlinear} transfer vibration energy to
an attached subsystem, where it is dissipated.
Energy-based control offers the same mechanism in software: within the
dissipativity and passivity frameworks~\citep{willems1972dissipative,
vanderschaft2000l2}, controlled-Lagrangian
methods~\citep{bloch2000controlled,woolsey2004controlled} and
interconnection-and-damping assignment~\citep{ortega2002interconnection} shape
both the stored energy and the interconnection structure of a closed loop.
A \emph{virtual mechanical system} (VMS) evolves an auxiliary mechanical model in
the controller and couples it to the plant through sensors and actuators, so that
the energy-transfer mechanism of a physical attachment is retained while its
parameters become freely tunable~\citep{preumont2018vibration,juchem2024virtual}.

This paper studies the simplest instance in which the interconnection itself is
the design variable: a host oscillator $(q,\dot q)$ coupled to a virtual
oscillator $(z,\dot z)$ through a constant skew-symmetric (gyroscopic) velocity
coupling $n$.
The interconnection preserves the total energy and only redistributes it.
In~\citet{juchem2026shaping} the closed loop was shown to evolve on invariant
two-tori whose projections on the host phase plane are Lissajous curves.
For a coprime modal ratio $\Op/\Om=\tau/\sigma$ these curves close, and they
pass through the host origin exactly when $\tau-\sigma\equiv2\pmod4$.
The host energy is then transferred \emph{completely} to the controller once
per common period, and a dissipative stage can harvest it.
This makes the coupling $n$ a tuning knob that trades depth of energy
exchange against speed.

Such an exchange geometry is a property of an \emph{ideal} skew-symmetric
interconnection.
Implementation breaks skew-symmetry: the force $n\dot z$ is applied by an
actuator, the reciprocal input $-n\dot q$ is computed from a measured, sampled and
delayed velocity, and gains, delays, zero-order holds and backlash all enter the
loop.
\citet{juchem2027restoring} observed in simulation that such imperfections cause
energy drift, unequal modal participation and loss of orbit closure, and proposed
a hierarchy of three corrective loops acting on the antisymmetric coupling, the
common coupling and the virtual stiffness.
Its analysis was windowed, it assumed boundedness and persistent excitation along
the closed-loop trajectory, and the frequency loop was only shown to be locally
convergent.
The present paper replaces that analysis with structural results and, as a
consequence, changes part of the design.

\paragraph*{Related work}
The loss of passivity of a virtual coupling under sampling is classical in
haptics~\citep{colgate1994passivity}; time-domain passivity
control~\citep{hannaford2002time} and discrete port-Hamiltonian
formulations~\citep{stramigioli2005sampled} restore passivity by dissipating the
energy generated by the implementation.
Here, however, the goal is not passivity of the interconnection but
\emph{conservation of a specific resonant geometry}, so energy has to be
redistributed among the modes rather than dissipated.
Delays in active vibration absorbers have been exploited deliberately, as in the
delayed resonator~\citep{olgac1994novel}.
The stability of linear retarded systems with a single delay can be decided
exactly through imaginary-axis crossings~\citep{cooke1986zeroes,
walton1987direct,olgac2002exact,michiels2014stability}, without Pad\'e
approximation.
Adaptive tuned absorbers retune an auxiliary oscillator under
mistuning~\citep{bonello2011adaptive}.
Phase locking of coupled oscillators~\citep{adler1946study,
pikovsky2001synchronization} is the natural language for the relative-phase
condition derived below.
On the analysis side we rely on input-to-state stability
(ISS)~\citep{sontag1989smooth,sontag2008input} and its preservation under
averaging~\citep{nesic2001input,teel2003unified} and
sampling~\citep{nesic2004framework}.

\paragraph*{Contributions}
\begin{enumerate}
\item \emph{Exact geometry.} A linear modal--polar change of coordinates turns
  the three defects of~\citet{juchem2027restoring} into coordinates: total
  energy, the log-amplitude ratio $e_A$, and the coprime relative phase
  $\psi=\sigma\theta_+-\tau\theta_-$
  (\cref{lem:modal,lem:perturbed}).
  Complete host absorption occurs if and only if
  $-\ln(\tau/\sigma)\le e_A\le0$ and $\psi$ lies in an explicit finite set,
  e.g.\ $\psi\equiv(\tau+\sigma)(2j+1)\pi/2$ for $e_A=0$
  (\cref{thm:absorption}).
  The $\tau-\sigma\equiv2\pmod4$ criterion of~\citet{juchem2026shaping} follows
  as the special case of the impulse response, and \emph{every} coprime pair
  becomes absorbing once the phase is controlled.
\item \emph{Exact fragility results.} For every positive loop delay the
  interconnection is exponentially unstable, with a closed-form count of the
  unstable roots; the fast mode is destabilised and the slow mode damped
  (\cref{thm:delay}).
  A zero-order hold is, to first order, an additional loop delay of one sampling
  period (\cref{prop:zoh}).
  Retuning the virtual stiffness tolerates only an $n^2/8$ relative coupling
  increase, while a trim of the common coupling has $O(1)$ authority
  (\cref{prop:authority}).
  This explains the non-convergence of the frequency loop
  of~\citet{juchem2027restoring} observed in \cref{sec:sim}.
\item \emph{Control and ISS.} Scale-invariant energy and amplitude laws, the
  latter with integral action, and a relative-phase locking law acting on the
  common coupling.
  Their first-order averaged dynamics are explicit and triangular
  (\cref{lem:averaged}) and ISS with respect to delay, measurement and gain
  errors (\cref{thm:iss-avg}).
  Semiglobal practical ISS of the actual sampled loop follows by averaging
  (\cref{thm:spiss}).
\end{enumerate}

\paragraph*{Notation}
$\RR_{>0}$ and $\NN$ denote the positive reals and positive integers;
$\lfloor\cdot\rfloor$ is the floor function; $\ii^2=-1$.
For $x\in\RR^4$, $|x|$ is the Euclidean norm, and $\|v\|_\infty$ is the
(essential) supremum norm of a signal.
A function $\gamma:\RR_{\ge0}\to\RR_{\ge0}$ is of class $\mathcal K$ if it is
continuous, zero at zero and strictly increasing; $\beta$ is of class
$\mathcal{KL}$ if $\beta(\cdot,t)\in\mathcal K$ and $\beta(r,\cdot)$ decreases
to zero.
For a $P_0$-periodic or almost-periodic function $f$,
$\avg{f}=\lim_{T\to\infty}T^{-1}\int_0^Tf\,\dd t$.

\paragraph*{Organisation}
\Cref{sec:model} introduces the model and the modal--polar coordinates.
\Cref{sec:tuning} recalls the resonant tuning framework
of~\citet{juchem2026shaping} and proves the absorption theorem.
\Cref{sec:mismatch} analyses delay, sampling and static mismatch.
\Cref{sec:control} presents the controller and its stability analysis,
\cref{sec:sim} the simulations, and \cref{sec:concl} concludes.
The longer proofs are in the Appendix.
